\SourceLeaf{436}%
\Needspace{7\baselineskip}\begin{center}13.\end{center}
\noindent\hspace*{1.2em}Jedem Punkt einer gekrümmten Linie auf der gekrümmten Fläche entsprechen in unserer Darstellung bereits zwei Punkte auf der Kugeloberfläche, nämlich der Punkt $\lambda$, der die Richtung des Linienelements angibt, und der Punkt $L$, der die Richtung der Normalen auf die Fläche angibt. Beide sind offenbar um $90^{0}$ voneinander entfernt. Bei unserer früheren Untersuchung (Art. $9$), in der [wir] vorausgesetzt hatten, dass die gekrümmte Linie in einer einzigen Ebene liegt, hatten wir noch \textls[100]{zwei} weitere Punkte auf der Kugeloberfläche, nämlich $\mathfrak{L}$, der die Richtung der Normalen auf die Ebene angab, und $\lambda'$, der die Richtung der Normalen auf das Kurvenelement in der Ebene angab. Dabei war $\mathfrak{L}$ ein fester Punkt, und $\lambda$, $\lambda'$ lagen stets auf einem einzigen Großkreis mit dem Pol $\mathfrak{L}$. Bei der Verallgemeinerung der Betrachtung wollen wir die Bezeichnungen $\mathfrak{L}$, $\lambda'$ beibehalten, müssen ihre Bedeutung aber von einem allgemeineren Standpunkt aus auffassen. Beim Durchlaufen der Kurve $s$ beschreiben die Punkte $L,\lambda$ auf der Hilfskugel ebenfalls Kurven, die im Allgemeinen keine Großkreise mehr sind. Man ziehe parallel zum Element der zweiten Kurve einen Radius der Hilfskugel bis zum Punkt $\lambda'$. Wir wählen stattdessen aber den gegenüberliegenden Punkt, wenn jener mehr als $90^{0}$ von $L$ entfernt ist. Im ersten Fall betrachten wir das Element bei $\lambda$ als positiv, im anderen als negativ. Schließlich sei $\mathfrak{L}$ der Punkt der Hilfskugel, der von $\lambda$ und $\lambda'$ zugleich um $90^{0}$ entfernt ist, und zwar so, dass $\lambda,\lambda',\mathfrak{L}$ in derselben Reihenfolge liegen wie $(1),(2),(3)$.

Die Koordinaten der vier Punkte auf der Hilfskugel in Bezug auf den Mittelpunkt sind
\[
\begin{array}{cccc}
L&X&Y&Z\\
\lambda&\xi&\eta&\zeta\\
\lambda'&\xi'&\eta'&\zeta'\\
\mathfrak{L}&\alpha&\beta&\gamma.
\end{array}
\]
Jeder dieser $4$ Punkte beschreibt also auf der Hilfskugel selbst eine Linie, deren Elemente wir mit $dL,d\lambda,d\lambda',d\mathfrak{L}$ bezeichnen wollen. Somit gilt
